% Options for packages loaded elsewhere
% Options for packages loaded elsewhere
\PassOptionsToPackage{unicode}{hyperref}
\PassOptionsToPackage{hyphens}{url}
\PassOptionsToPackage{dvipsnames,svgnames,x11names}{xcolor}
%
\documentclass[
  letterpaper,
  DIV=11,
  numbers=noendperiod]{scrreprt}
\usepackage{xcolor}
\usepackage{amsmath,amssymb}
\setcounter{secnumdepth}{5}
\usepackage{iftex}
\ifPDFTeX
  \usepackage[T1]{fontenc}
  \usepackage[utf8]{inputenc}
  \usepackage{textcomp} % provide euro and other symbols
\else % if luatex or xetex
  \usepackage{unicode-math} % this also loads fontspec
  \defaultfontfeatures{Scale=MatchLowercase}
  \defaultfontfeatures[\rmfamily]{Ligatures=TeX,Scale=1}
\fi
\usepackage{lmodern}
\ifPDFTeX\else
  % xetex/luatex font selection
\fi
% Use upquote if available, for straight quotes in verbatim environments
\IfFileExists{upquote.sty}{\usepackage{upquote}}{}
\IfFileExists{microtype.sty}{% use microtype if available
  \usepackage[]{microtype}
  \UseMicrotypeSet[protrusion]{basicmath} % disable protrusion for tt fonts
}{}
\makeatletter
\@ifundefined{KOMAClassName}{% if non-KOMA class
  \IfFileExists{parskip.sty}{%
    \usepackage{parskip}
  }{% else
    \setlength{\parindent}{0pt}
    \setlength{\parskip}{6pt plus 2pt minus 1pt}}
}{% if KOMA class
  \KOMAoptions{parskip=half}}
\makeatother
% Make \paragraph and \subparagraph free-standing
\makeatletter
\ifx\paragraph\undefined\else
  \let\oldparagraph\paragraph
  \renewcommand{\paragraph}{
    \@ifstar
      \xxxParagraphStar
      \xxxParagraphNoStar
  }
  \newcommand{\xxxParagraphStar}[1]{\oldparagraph*{#1}\mbox{}}
  \newcommand{\xxxParagraphNoStar}[1]{\oldparagraph{#1}\mbox{}}
\fi
\ifx\subparagraph\undefined\else
  \let\oldsubparagraph\subparagraph
  \renewcommand{\subparagraph}{
    \@ifstar
      \xxxSubParagraphStar
      \xxxSubParagraphNoStar
  }
  \newcommand{\xxxSubParagraphStar}[1]{\oldsubparagraph*{#1}\mbox{}}
  \newcommand{\xxxSubParagraphNoStar}[1]{\oldsubparagraph{#1}\mbox{}}
\fi
\makeatother


\usepackage{longtable,booktabs,array}
\usepackage{calc} % for calculating minipage widths
% Correct order of tables after \paragraph or \subparagraph
\usepackage{etoolbox}
\makeatletter
\patchcmd\longtable{\par}{\if@noskipsec\mbox{}\fi\par}{}{}
\makeatother
% Allow footnotes in longtable head/foot
\IfFileExists{footnotehyper.sty}{\usepackage{footnotehyper}}{\usepackage{footnote}}
\makesavenoteenv{longtable}
\usepackage{graphicx}
\makeatletter
\newsavebox\pandoc@box
\newcommand*\pandocbounded[1]{% scales image to fit in text height/width
  \sbox\pandoc@box{#1}%
  \Gscale@div\@tempa{\textheight}{\dimexpr\ht\pandoc@box+\dp\pandoc@box\relax}%
  \Gscale@div\@tempb{\linewidth}{\wd\pandoc@box}%
  \ifdim\@tempb\p@<\@tempa\p@\let\@tempa\@tempb\fi% select the smaller of both
  \ifdim\@tempa\p@<\p@\scalebox{\@tempa}{\usebox\pandoc@box}%
  \else\usebox{\pandoc@box}%
  \fi%
}
% Set default figure placement to htbp
\def\fps@figure{htbp}
\makeatother





\setlength{\emergencystretch}{3em} % prevent overfull lines

\providecommand{\tightlist}{%
  \setlength{\itemsep}{0pt}\setlength{\parskip}{0pt}}



 


\KOMAoption{captions}{tableheading}
\makeatletter
\@ifpackageloaded{bookmark}{}{\usepackage{bookmark}}
\makeatother
\makeatletter
\@ifpackageloaded{caption}{}{\usepackage{caption}}
\AtBeginDocument{%
\ifdefined\contentsname
  \renewcommand*\contentsname{Table of contents}
\else
  \newcommand\contentsname{Table of contents}
\fi
\ifdefined\listfigurename
  \renewcommand*\listfigurename{List of Figures}
\else
  \newcommand\listfigurename{List of Figures}
\fi
\ifdefined\listtablename
  \renewcommand*\listtablename{List of Tables}
\else
  \newcommand\listtablename{List of Tables}
\fi
\ifdefined\figurename
  \renewcommand*\figurename{Figure}
\else
  \newcommand\figurename{Figure}
\fi
\ifdefined\tablename
  \renewcommand*\tablename{Table}
\else
  \newcommand\tablename{Table}
\fi
}
\@ifpackageloaded{float}{}{\usepackage{float}}
\floatstyle{ruled}
\@ifundefined{c@chapter}{\newfloat{codelisting}{h}{lop}}{\newfloat{codelisting}{h}{lop}[chapter]}
\floatname{codelisting}{Listing}
\newcommand*\listoflistings{\listof{codelisting}{List of Listings}}
\makeatother
\makeatletter
\makeatother
\makeatletter
\@ifpackageloaded{caption}{}{\usepackage{caption}}
\@ifpackageloaded{subcaption}{}{\usepackage{subcaption}}
\makeatother
\usepackage{bookmark}
\IfFileExists{xurl.sty}{\usepackage{xurl}}{} % add URL line breaks if available
\urlstyle{same}
\hypersetup{
  pdftitle={Sociology of Water Management},
  pdfauthor={A. Buday},
  colorlinks=true,
  linkcolor={blue},
  filecolor={Maroon},
  citecolor={Blue},
  urlcolor={Blue},
  pdfcreator={LaTeX via pandoc}}


\title{Sociology of Water Management}
\author{A. Buday}
\date{2025-12-30}
\begin{document}
\maketitle

\renewcommand*\contentsname{Table of contents}
{
\hypersetup{linkcolor=}
\setcounter{tocdepth}{2}
\tableofcontents
}

\bookmarksetup{startatroot}

\chapter{Sociology of Water
Management}\label{sociology-of-water-management}

The waters of the Laurentian Great Lakes are a special treasure in a
thirsty world. More than 30 million people in eight U.S. states and one
Canadian province rely on this system for drinking water, food
production, transportation, recreation, and cultural identity, making
management of this precious resource an inherently collaborative
endeavor.

In recent years, there has been increasing interest in integrating
social science with ecological research to develop management approaches
that account for the human dimensions of natural resource management.
From fisheries to farming and all things in between, human activities
shape Great Lakes ecosystems in complex and far-reaching ways.

This chapter examines these dynamics, including legacy pollution from
industrialization and urbanization in the basin, as well as present-day
management challenges associated with modern sources of pollution.

\bookmarksetup{startatroot}

\chapter{From Data To Statistics}\label{from-data-to-statistics}

Learning Objectives

\begin{itemize}
\tightlist
\item
  Organize raw data using frequency distributions
\item
  Describe the shape, center, and variation of a variable
\item
  Explain what descriptive statistics tell us about survey responses
\end{itemize}

I was a horrible statistics student. I approached the subject with a
survivor mentality, clawing my way through assignments and exams and
immediately forgetting all related content due to immense anxiety. I
took my B in graduate social statistics (basically failing in graduate
terms) and gladly moved on to complete a qualitative dissertation.

Imagine my surprise when I landed my first job and was immediately asked
to lead projects requiring survey research. I did what any young
graduate with a mountain of student loan debt does -- jumped right in!

When I first encountered John Kranzler's (2017) book, \emph{Statistics
for the Terrified}, I felt seen. Kranzler had a way of describing my
statistics paralysis in hilarious prose and breaking down basic
terminology in a way I found accessible -- and even fun. My own take on
this topic therefore follows his excellent model.

If you are feeling stressed about talking about statistics, I invite you
to take a moment to watch the short video:
\href{https://youtu.be/iymgY56mgCo?si=4t7mUiiaIKp-mhCh}{What is
Statistics Anxiety?}

Stats anxiety is problematic because it blocks us from learning,
prevents us from asking questions, and generally snowballs into an
unproductive physical and mental state. Instead of focusing on your
performance goals and the fear that you won't reach them, try focusing
on the learning process as we move forward, and acknowledge how
struggling with new material can contribute to your academic and
professional growth.

We will ease into our data explorations by reviewing the structure of
data in raw form and tracing the process of transitioning data from
unordered numbers to interpretable distributions. This introduction will
have relevance beyond your West MI Watersheds analysis, laying a
foundation of data literacy you will build on throughout your
professional life.

\section{Data Structure}\label{data-structure}

Let's begin our journey into data analytics by clarifying two important
terms: data and statistics.

In quantitative research, data are pieces of information represented as
numbers. These numbers may describe things like behaviors, experiences,
opinions, or characteristics of people. This contrasts with qualitative
data, which are expressed as text, images, or other non-numeric forms
(Sheppard, 2020).

Statistics refers to the set of tools we use to organize, summarize, and
display quantitative data in ways that make patterns easier to see and
interpret (Kranzler, 2017). We'll begin with the most basic statistical
task of all: turning raw data into something we can understand.

Before we can describe or analyze data, we need to understand what we
are actually looking at when we open a dataset in RStudio. At first
glance, a data file looks like a large grid of numbers that doesn't mean
very much at all. In reality, data files follow a conventional and
meaningful structure:

\begin{itemize}
\tightlist
\item
  Each row represents a case. In survey data, each row corresponds to
  one survey participant.
\item
  Each column represents a variable, or a characteristic measured by the
  survey (such as age, frequency of septic tank pumping, or perceptions
  of water quality).
\item
  Each cell contains a value, which is the recorded response for one
  person on one variable (Çetinkaya-Rundel \& Hardin, 2024).
\end{itemize}

Understanding this structure helps you keep track of which information
lives where in your dataset. If you are interested in how people
answered a particular survey question, you will look down a column. If
you want to understand one participant's full set of responses, you will
look across a row.

However, simply seeing numbers in a dataset does not mean those numbers
can be analyzed immediately. Raw data often include values that serve
special purposes, such as indicating a skipped question, a ``don't
know'' response, or a response that does not apply to a particular
person (Arthur \& Clark, 2021). These values are part of the dataset's
structure, but they require interpretation before analysis.

Another critical concept to understand is the difference between a
missing value and a value of zero (Zimmer et al., 2024).

In RStudio, missing data are represented as \texttt{NA}, which stands
for ``not available.'' An \texttt{NA} means that no valid response was
recorded for that cell. That might occur because a respondent skipped a
question, was not eligible to answer it, or recorded a response that was
invalid (such as selecting two answers to the same question or penciling
in a response off the given answer scale).

A value of \texttt{0}, by contrast, is not missing. A zero is a recorded
value that has meaning defined in the data documentation. For example, a
value of \texttt{0} might indicate that a respondent did not engage in a
particular activity, used a practice zero times, or selected a specific
response option that was coded \texttt{0}.

This distinction matters because R treats \texttt{NA} and \texttt{0}
very differently. Missing values are excluded from most calculations by
default, while zero values are included. If you accidentally treat a
zero as missing, or a missing value as a zero, you can seriously distort
your results.

For these reason, one of your first responsibilities when working with a
dataset is to use the codebook to determine which values represent real
responses and which represent missing or non-applicable information.
Only then can you begin summarizing your data with confidence.

Glossary Terms: Statistics: Techniques and procedures for working with
quantitative data. Data: A collection of information. Raw data: Data in
its original, unorganized form before it has been summarized or cleaned.
Missing value: Data entries indicating that no valid response is
available; coded as NA in R.

\section{Data Documentation}\label{data-documentation}

Once you understand the basic structure of a dataset -- rows as people,
columns as variables, and cells as values -- the next step is learning
how to interpret what those values actually mean. This is where data
documentation, and especially the codebook, becomes essential.

Without documentation a dataset is just a jumble of numbers. With
documentation, those numbers become meaningful information about people,
behaviors, and perceptions.

A codebook (sometimes called a data dictionary) serves as the
translation key between human language and machine-readable data (Arthur
\& Clark, 2021).

At minimum, a codebook typically tells you:

\begin{itemize}
\tightlist
\item
  The variable name used in the dataset
\item
  A description of what the variable measures (the survey question)
\item
  The value label assigned to each response option
\item
  Any special codes used for ``don't know'' or inapplicable responses
  (Zimmer et al., 2024)
\end{itemize}

For example, a survey question asking about the frequency of septic tank
pumping may be stored as a variable named SEPTICPUMP, with values coded
as:

\begin{itemize}
\tightlist
\item
  0 = Never
\item
  1 = More than 10 years
\item
  2 = Every 6-10 years
\item
  3 = Every 3-5 years
\item
  998 = I don't have a septic system
\end{itemize}

Looking only at the dataset, you would see numbers like
\texttt{0,\ 1,\ 2,\ 3,\ 998}. The codebook tells you how to interpret
those numbers correctly.

It may be tempting to jump straight into calculating averages, creating
graphs, or running models as soon as you load your data in R. However,
doing so without consulting the codebook can lead to serious errors
(Sheppard, 2020). Because of this, you should always examine the
codebook before summarizing a variable and be familiar with the items as
they were presented to respondents in the questionnaire, or survey
instrument, itself (Zimmer et al., 2024).

Most surveys include responses that are not part of an ordered scale,
such as ``don't know'' or ``not applicable.'' These responses are often
assigned out-of-range values (such as 9, 99, or 998) so that they are
not confused with substantive responses (Arther \& Clark, 2021; Zimmer
et al., 2024). These numeric codes may need to be converted to missing
data (NA) before calculating means or medians, or they will distort the
reported values. Considering our perceptions of water quality variables,
above, if I calculated an average for these items without reassigning
out-of-range values to missing data, I might end up with an average
score of 6 on an item whose substantive range is 1-4. Oops!

Beyond interpreting values, codebooks also guide common data management
tasks, such as deciding when to create dummy variables (0/1 indicators),
identifying variables that need reverse coding, understanding how
multi-item concepts are measured, and determining which variables are
comparable across survey locations. These transformations are a routine
part of quantitative data preparation (Arthur \& Clark, 2021).

Using documentation may feel tedious or confusing at first, but it is
one of the most important habits you will develop as a data analyst.
Careful attention to codebooks prevents misinterpretation, supports
transparency, and ensures that your findings are defensible. Before
moving on to describing distributions or calculating statistics, make it
routine to ask:

\begin{itemize}
\tightlist
\item
  What does this variable measure?
\item
  What do the values represent?
\item
  Are any responses coded in a special way?
\item
  Do I need to clean or recode this variable before summarizing it?
\end{itemize}

These questions will guide you toward more accurate analysis.

Glossary Terms: Data documentation: The collection of materials, such as
codebooks, questionnaires, and technical notes, that describe how data
were collected, coded, and structured. Codebook: A document that
explains how survey responses were recorded and stored in a dataset.
Questionnaire: The survey instrument presented to respondents. Value
label: Descriptions that explain what specific numeric values represent
(e.g., 1 = ``Never,'' 3 = ``Every 3--5 years''). Out-of-range value: A
numeric code (such as 9, 99, or 998) used to represent responses like
``don't know'' or ``not applicable'' that fall outside the substantive
response scale. Data management: The process of preparing data for
analysis, including cleaning, recoding, transforming variables, and
handling missing or special values. Dummy variables: A binary variable
coded as 0 or 1 to indicate the absence or presence of a specific
characteristic or response. Reverse coding: A data management technique
that flips the direction of a scale so higher values consistently
represent the same underlying concept.

\section{Frequency distributions}\label{frequency-distributions}

We've talked all around it so now let's get down to business -- it's
time to look at our survey data. When I begin any data exploration, the
first thing I want to do is get familiar with the distribution of
responses on each item in my analysis. This gives me an understanding of
how people responded to each survey question and whether or not an
individual item is a good candidate for use in more advanced analyses.

As we've reviewed, data in their raw form appear as unordered sets of
scores that don't tell us very much. A frequency distribution organizes
data by summarizing how often a specific score was obtained (Kranzler,
2017).

For example, take a look at these quiz scores from my SOC 105: Social
Problems course. What can we know about this group of quiz scores at
first glance?

Not much! If you scored 12/15 points, what would that tell you? Did you
do better or worse than the average student? To answer these questions,
we need to organize our quiz score data by creating a frequency
distribution.

The first step is to identify the highest and lowest score in the set.

Next, I'll make a list of all possible scores, from lowest to highest.

Then, I'll go through the scores, one-by-one, and make a tally mark each
time a score occurs.

Finally, I'll count the number of tallies to find the frequency (f) with
which each score appeared in the dataset.

What can we know about this group of quiz scores now? More than before!
The scores range from 6-15. The score obtained most often was 13. Most
scores were between 10-14, so your score (12) was close to the class
average.

Displaying frequency distributions graphically can be helpful for
describing the shape of the distribution of scores on a given variable
(Agresti, 2018).

Observing the shape of the distribution is important for determining the
extent to which responses on an item varied evenly, diverged sharply, or
grouped around a single answer option.

A bell-shaped distribution indicates that most scores fall near a
central value. This is also called a ``normal'' curve. The distribution
is symmetrical.

In Contrast, a U-shaped distribution indicates that most scores fall
near the highest and lowest possible values, meaning that the variable
is polarized.

Additionally, a {[}pb\_glossary id=``595''{]}skewed
distribution{[}/pb\_glossary{]} has values that cluster around a common
score. The distribution may be negatively (right) skewed, with most
scores clustered at the high end of the distribution and the tail of the
distribution longer on the left side than the right. Or, it may be
positively (left) skewed, with most scores clustered at the low end of
the distribution and the tail of the distribution longer on the right
side than the left.

The shape of the distribution matters because it affects what we can
reasonably do with the data next. Variables that follow a normal
distribution are especially analytically rich. When responses are
distributed symmetrically around a central value, common summary
statistics provide a meaningful description of a ``typical'' case. Many
statistical techniques assume a normal distribution, and highly skewed
distributions make it difficult to assess differences between groups.
All distributions tell us something meaningful about people's responses,
but their analytic power and suitability for further analysis can vary.

Glossary Terms: Bell-shaped distribution: a distribution in which most
scores fall near a central value, producing a symmetrical shape.
Frequency distribution: a summary that describes how often a specific
score appears in a dataset. Skewed distribution: a distribution in which
most scores fall near the high or low end of the values possible,
producing a long tail on one end.

\section{Summary}\label{summary}

This chapter focused on the transition from raw data to meaningful
statistical insight. Rather than rushing into calculations, we
emphasized the importance of slowing down to understand how data are
structured, documented, and interpreted before analysis begins. Doing so
helps ensure that the numbers we work with accurately reflect people's
responses and experiences.

We explored how organizing data into frequency distributions transforms
unordered values into interpretable patterns. Looking at how responses
are distributed allows us to see where answers cluster, how much they
vary, and whether a variable is likely to support more advanced
analysis. Just as importantly, we highlighted that not all data are
immediately ready for analysis. Values may require interpretation,
cleaning, or recoding before they can be summarized responsibly.

Together, these steps establish a foundation for descriptive statistics:
understanding what the data represent, how responses are distributed,
and why those patterns matter. Moving forward, we'll build on this
foundation by introducing levels of measurement and work on formally
describing the center and spread of distributions. Keep calm and
statistics on!

\section{Reflection Questions}\label{reflection-questions}

\begin{enumerate}
\def\labelenumi{\arabic{enumi}.}
\item
  \textbf{Survey datasets often include values that represent different
  kinds of information, such as substantive responses, missing data, or
  special codes like ``don't know.''} How might misunderstanding these
  distinctions affect the conclusions a researcher draws from a dataset?
\item
  \textbf{Frequency distributions are often the first step in
  quantitative analysis.} How does examining the distribution of
  responses help researchers assess whether a survey question
  meaningfully captures variation in people's experiences or opinions?
\item
  \textbf{The shape of a distribution can reveal clustering,
  polarization, or imbalance in responses.} Why is understanding these
  patterns important before applying more advanced statistical
  techniques?
\item
  \textbf{Early decisions about how data are organized, cleaned, and
  summarized shape everything that follows in an analysis.} How might
  careless early steps limit the usefulness or credibility of later
  findings?
\end{enumerate}

\section{References}\label{references}

Agresti, A. (2018). Statistical methods for the social sciences (5th
ed.). Pearson Education.

Arthur, M. M. L., \& Clark, R. (2021). Social data analysis: Qualitative
and quantitative approaches. Open Educational Resource.

Çetinkaya-Rundel, M. and Hardin, J. (2024). Introduction to Modern
Statistics, 2e. OpenIntro. CC-SA 3.0
(https://openintro-ims.netlify.app/).

Kranzler, J.H. (2017). Statistics for the Terrified (6th ed.). Roman \&
Littlefield.

Sheppard, V. (2020). Research Methods for the Social Sciences: An
Introduction. Pressbooks. CC-BY-NC-SA 4.0.
(https://pressbooks.bccampus.ca/jibcresearchmethods/).

Zimmer, S. A., Powell, R. J., \& Velásquez, I. C. (2024). Exploring
Complex Survey Data Analysis Using R: A Tidy Introduction with \{srvyr\}
and \{survey\}. Chapman \& Hall: CRC Press.
(https://tidy-survey-r.github.io/tidy-survey-book/).




\end{document}
